% --- Archon physics formalization source begin ---
% archon:physics
% archon:covers PhyXMiniProblems/problem_phyx_mini_0269.lean
% archon:source-report reports/phyx_mini/problem_phyx_mini_0269.source.json

\chapter{Physics problem phyx\_mini\_0269}
\label{ch:PhyXMiniProblems_problem_phyx_mini_0269}

\paragraph{Problem source.}
\#\# Physical scenario

In figure, a metal plate is mounted on an axle through its center of mass. A spring with \$k = 2000\textbackslash{} N/m\$ connects a wall with a point on the rim a distance \$r = 2.5\textbackslash{} cm\$ from the center of mass. Initially the spring is at its rest length. If the plate is rotated by \$7\textasciicircum{}\textbackslash{}circ\$ and released, it rotates about the axle in SHM, with its angular position given by figure. The horizontal axis scale is set by \$t\_s = 20\textbackslash{} ms\$.

\#\# Question

What is the rotational inertia of the plate about its center of mass?

\#\# Answer choices

- A: A: \$0.8\textbackslash{}times 10\textasciicircum{}\{-5\}\textbackslash{}\textbackslash{} kg \textbackslash{}cdot m\textasciicircum{}2\$

- B: B: \$1.0\textbackslash{}times 10\textasciicircum{}\{-5\}\textbackslash{}\textbackslash{} kg \textbackslash{}cdot m\textasciicircum{}2\$

- C: C: \$1.3\textbackslash{}times 10\textasciicircum{}\{-5\}\textbackslash{}\textbackslash{} kg \textbackslash{}cdot m\textasciicircum{}2\$

- D: D: \$1.6\textbackslash{}times 10\textasciicircum{}\{-5\}\textbackslash{}\textbackslash{} kg \textbackslash{}cdot m\textasciicircum{}2\$

\#\# Figure caption (auxiliary; use the image as primary evidence)

The image contains two main parts:

1. **Left Side:**
   - An illustration of a wall-mounted setup with a spring.
   - The spring is horizontally attached to a gear-like object or a jagged circular shape.
   - The object with the center dot appears to represent a point of rotation or pivot.
   - The variable "r" is marked with a vertical arrow indicating the distance from the rotation point to the edge of the object.

2. **Right Side:**
   - A graph plotting angular displacement (θ in degrees) against time (t in milliseconds).
   - The θ-axis ranges from -8 to 8 degrees, and the t-axis is marked in milliseconds with divisions.
   - The graph shows a sinusoidal wave pattern, indicating oscillatory motion.
   - The horizontal axis crosses the graph at zero degrees.
   - The point \textbackslash{}( t\_s \textbackslash{}) is labeled on the time axis, indicating a specific time point in the oscillation.

Overall, the illustration and graph depict a system experiencing harmonic or oscillatory motion.

\#\# Dataset metadata

Resonance and Harmonics; Physical Model Grounding Reasoning, Multi-Formula Reasoning

\paragraph{Recorded answer/context.}
C: C: \$1.3\textbackslash{}times 10\textasciicircum{}\{-5\}\textbackslash{}\textbackslash{} kg \textbackslash{}cdot m\textasciicircum{}2\$

\paragraph{Figure/image path.}
/root/proposal\_for\_physic/science-mango/phyx\_mini\_run/phyx\_data/test\_image/269.png

\paragraph{Formalization target.}
create a compiling Lean file with sorry bodies at `PhyXMiniProblems/problem\_phyx\_mini\_0269.lean`. The Lean declarations must preserve the physical quantities, dimensions or dimensional roles, figure labels, governing-law hypotheses, and final relation expressed by this problem.
Use Mathlib/Physlib names found through LeanExplore where available. If a physics API is missing, introduce faithful local abstractions rather than scalar placeholder aliases.

\begin{theorem}[Physics formalization target]
\label{thm:physics:phyx_mini_0269:target}
\lean{PhyXMiniProblems.ProblemPhyXMini0269.problem_phyx_mini_0269}
\uses{def:physics:phyx-mini-0269:phyxminiproblems-problemphyxmini0269-lengthinmeters, def:physics:phyx-mini-0269:phyxminiproblems-problemphyxmini0269-timeinseconds, def:physics:phyx-mini-0269:phyxminiproblems-problemphyxmini0269-springconstantinnewtonspermeter, def:physics:phyx-mini-0269:phyxminiproblems-problemphyxmini0269-rotationalinertiainkilogrammeterssquared, def:physics:phyx-mini-0269:phyxminiproblems-problemphyxmini0269-platespringsetup, def:physics:phyx-mini-0269:phyxminiproblems-problemphyxmini0269-matchesverbalscenario, def:physics:phyx-mini-0269:phyxminiproblems-problemphyxmini0269-matchesproblemandfigurereadouts, def:physics:phyx-mini-0269:phyxminiproblems-problemphyxmini0269-hasphysicalplatespringparameters, def:physics:phyx-mini-0269:phyxminiproblems-problemphyxmini0269-satisfieslinearizedplatespringlaws, def:physics:phyx-mini-0269:phyxminiproblems-problemphyxmini0269-satisfiesidealangularshmlaw, def:physics:phyx-mini-0269:phyxminiproblems-problemphyxmini0269-isuniqueclosestanswer}
The assigned autoformalize agent should translate this physics problem into Lean declarations in the covered file, with theorem and lemma proof bodies written as `by sorry`.
\end{theorem}
\begin{proof}
This is an autoformalization task, not a proof task. Produce statements that can later be proved by the physics prover without weakening the physical model.
\end{proof}

% --- Archon declaration topology begin ---
\section{Declaration topology}
\begin{definition}[Length Quantity]
  \label{def:physics:phyx-mini-0269:phyxminiproblems-problemphyxmini0269-lengthquantity}
  \lean{PhyXMiniProblems.ProblemPhyXMini0269.LengthQuantity}
  A nonnegative physical length
\end{definition}
\begin{proof}
  This is the declared model definition of Length Quantity.
\end{proof}
\begin{definition}[Signed Length Quantity]
  \label{def:physics:phyx-mini-0269:phyxminiproblems-problemphyxmini0269-signedlengthquantity}
  \lean{PhyXMiniProblems.ProblemPhyXMini0269.SignedLengthQuantity}
  A signed extension along the spring axis
\end{definition}
\begin{proof}
  This is the declared model definition of Signed Length Quantity.
\end{proof}
\begin{definition}[Time Quantity]
  \label{def:physics:phyx-mini-0269:phyxminiproblems-problemphyxmini0269-timequantity}
  \lean{PhyXMiniProblems.ProblemPhyXMini0269.TimeQuantity}
  A nonnegative physical duration
\end{definition}
\begin{proof}
  This is the declared model definition of Time Quantity.
\end{proof}
\begin{definition}[Spring Constant Quantity]
  \label{def:physics:phyx-mini-0269:phyxminiproblems-problemphyxmini0269-springconstantquantity}
  \lean{PhyXMiniProblems.ProblemPhyXMini0269.SpringConstantQuantity}
  A spring force constant has dimension force per length, or mass/time²
\end{definition}
\begin{proof}
  This is the declared model definition of Spring Constant Quantity.
\end{proof}
\begin{definition}[Signed Force Quantity]
  \label{def:physics:phyx-mini-0269:phyxminiproblems-problemphyxmini0269-signedforcequantity}
  \lean{PhyXMiniProblems.ProblemPhyXMini0269.SignedForceQuantity}
  A signed force component along the tangent at the rim attachment
\end{definition}
\begin{proof}
  This is the declared model definition of Signed Force Quantity.
\end{proof}
\begin{definition}[Rotational Inertia Quantity]
  \label{def:physics:phyx-mini-0269:phyxminiproblems-problemphyxmini0269-rotationalinertiaquantity}
  \lean{PhyXMiniProblems.ProblemPhyXMini0269.RotationalInertiaQuantity}
  Rotational inertia has dimension mass times length squared
\end{definition}
\begin{proof}
  This is the declared model definition of Rotational Inertia Quantity.
\end{proof}
\begin{definition}[Signed Torque Quantity]
  \label{def:physics:phyx-mini-0269:phyxminiproblems-problemphyxmini0269-signedtorquequantity}
  \lean{PhyXMiniProblems.ProblemPhyXMini0269.SignedTorqueQuantity}
  A signed torque about the axle through the plate's center of mass
\end{definition}
\begin{proof}
  This is the declared model definition of Signed Torque Quantity.
\end{proof}
\begin{definition}[Restoring Coefficient Quantity]
  \label{def:physics:phyx-mini-0269:phyxminiproblems-problemphyxmini0269-restoringcoefficientquantity}
  \lean{PhyXMiniProblems.ProblemPhyXMini0269.RestoringCoefficientQuantity}
  A nonnegative restoring-torque coefficient per dimensionless radian
\end{definition}
\begin{proof}
  This is the declared model definition of Restoring Coefficient Quantity.
\end{proof}
\begin{definition}[Angular Velocity Quantity]
  \label{def:physics:phyx-mini-0269:phyxminiproblems-problemphyxmini0269-angularvelocityquantity}
  \lean{PhyXMiniProblems.ProblemPhyXMini0269.AngularVelocityQuantity}
  Angular velocity has inverse-time dimension because radians are dimensionless
\end{definition}
\begin{proof}
  This is the declared model definition of Angular Velocity Quantity.
\end{proof}
\begin{definition}[Angular Acceleration Quantity]
  \label{def:physics:phyx-mini-0269:phyxminiproblems-problemphyxmini0269-angularaccelerationquantity}
  \lean{PhyXMiniProblems.ProblemPhyXMini0269.AngularAccelerationQuantity}
  Angular acceleration has inverse-time-squared dimension
\end{definition}
\begin{proof}
  This is the declared model definition of Angular Acceleration Quantity.
\end{proof}
\begin{definition}[Angular Frequency Quantity]
  \label{def:physics:phyx-mini-0269:phyxminiproblems-problemphyxmini0269-angularfrequencyquantity}
  \lean{PhyXMiniProblems.ProblemPhyXMini0269.AngularFrequencyQuantity}
  A nonnegative angular frequency
\end{definition}
\begin{proof}
  This is the declared model definition of Angular Frequency Quantity.
\end{proof}
\begin{definition}[nonnegative SIReadout]
  \label{def:physics:phyx-mini-0269:phyxminiproblems-problemphyxmini0269-nonnegativesireadout}
  \lean{PhyXMiniProblems.ProblemPhyXMini0269.nonnegativeSIReadout}
  Read a nonnegative dimensionful quantity in coherent SI units
\end{definition}
\begin{proof}
  This is the declared model definition of nonnegative SIReadout.
\end{proof}
\begin{definition}[signed SIReadout]
  \label{def:physics:phyx-mini-0269:phyxminiproblems-problemphyxmini0269-signedsireadout}
  \lean{PhyXMiniProblems.ProblemPhyXMini0269.signedSIReadout}
  Read a signed dimensionful quantity in coherent SI units
\end{definition}
\begin{proof}
  This is the declared model definition of signed SIReadout.
\end{proof}
\begin{definition}[length In Meters]
  \label{def:physics:phyx-mini-0269:phyxminiproblems-problemphyxmini0269-lengthinmeters}
  \lean{PhyXMiniProblems.ProblemPhyXMini0269.lengthInMeters}
  \uses{def:physics:phyx-mini-0269:phyxminiproblems-problemphyxmini0269-lengthquantity, def:physics:phyx-mini-0269:phyxminiproblems-problemphyxmini0269-nonnegativesireadout}
  Meter readout of a physical length
\end{definition}
\begin{proof}
  This is the declared model definition of length In Meters.
\end{proof}
\begin{definition}[signed Length In Meters]
  \label{def:physics:phyx-mini-0269:phyxminiproblems-problemphyxmini0269-signedlengthinmeters}
  \lean{PhyXMiniProblems.ProblemPhyXMini0269.signedLengthInMeters}
  \uses{def:physics:phyx-mini-0269:phyxminiproblems-problemphyxmini0269-signedlengthquantity, def:physics:phyx-mini-0269:phyxminiproblems-problemphyxmini0269-signedsireadout}
  Meter readout of a signed spring extension
\end{definition}
\begin{proof}
  This is the declared model definition of signed Length In Meters.
\end{proof}
\begin{definition}[time In Seconds]
  \label{def:physics:phyx-mini-0269:phyxminiproblems-problemphyxmini0269-timeinseconds}
  \lean{PhyXMiniProblems.ProblemPhyXMini0269.timeInSeconds}
  \uses{def:physics:phyx-mini-0269:phyxminiproblems-problemphyxmini0269-timequantity, def:physics:phyx-mini-0269:phyxminiproblems-problemphyxmini0269-nonnegativesireadout}
  Second readout of a physical duration
\end{definition}
\begin{proof}
  This is the declared model definition of time In Seconds.
\end{proof}
\begin{definition}[spring Constant In Newtons Per Meter]
  \label{def:physics:phyx-mini-0269:phyxminiproblems-problemphyxmini0269-springconstantinnewtonspermeter}
  \lean{PhyXMiniProblems.ProblemPhyXMini0269.springConstantInNewtonsPerMeter}
  \uses{def:physics:phyx-mini-0269:phyxminiproblems-problemphyxmini0269-springconstantquantity, def:physics:phyx-mini-0269:phyxminiproblems-problemphyxmini0269-nonnegativesireadout}
  SI readout of a spring constant, numerically in newtons per meter
\end{definition}
\begin{proof}
  This is the declared model definition of spring Constant In Newtons Per Meter.
\end{proof}
\begin{definition}[signed Force In Newtons]
  \label{def:physics:phyx-mini-0269:phyxminiproblems-problemphyxmini0269-signedforceinnewtons}
  \lean{PhyXMiniProblems.ProblemPhyXMini0269.signedForceInNewtons}
  \uses{def:physics:phyx-mini-0269:phyxminiproblems-problemphyxmini0269-signedforcequantity, def:physics:phyx-mini-0269:phyxminiproblems-problemphyxmini0269-signedsireadout}
  Newton readout of a signed tangential force
\end{definition}
\begin{proof}
  This is the declared model definition of signed Force In Newtons.
\end{proof}
\begin{definition}[rotational Inertia In Kilogram Meters Squared]
  \label{def:physics:phyx-mini-0269:phyxminiproblems-problemphyxmini0269-rotationalinertiainkilogrammeterssquared}
  \lean{PhyXMiniProblems.ProblemPhyXMini0269.rotationalInertiaInKilogramMetersSquared}
  \uses{def:physics:phyx-mini-0269:phyxminiproblems-problemphyxmini0269-rotationalinertiaquantity, def:physics:phyx-mini-0269:phyxminiproblems-problemphyxmini0269-nonnegativesireadout}
  Kilogram-meter-squared readout of rotational inertia
\end{definition}
\begin{proof}
  This is the declared model definition of rotational Inertia In Kilogram Meters Squared.
\end{proof}
\begin{definition}[signed Torque In Newton Meters]
  \label{def:physics:phyx-mini-0269:phyxminiproblems-problemphyxmini0269-signedtorqueinnewtonmeters}
  \lean{PhyXMiniProblems.ProblemPhyXMini0269.signedTorqueInNewtonMeters}
  \uses{def:physics:phyx-mini-0269:phyxminiproblems-problemphyxmini0269-signedtorquequantity, def:physics:phyx-mini-0269:phyxminiproblems-problemphyxmini0269-signedsireadout}
  Newton-meter readout of a signed torque
\end{definition}
\begin{proof}
  This is the declared model definition of signed Torque In Newton Meters.
\end{proof}
\begin{definition}[restoring Coefficient In Newton Meters]
  \label{def:physics:phyx-mini-0269:phyxminiproblems-problemphyxmini0269-restoringcoefficientinnewtonmeters}
  \lean{PhyXMiniProblems.ProblemPhyXMini0269.restoringCoefficientInNewtonMeters}
  \uses{def:physics:phyx-mini-0269:phyxminiproblems-problemphyxmini0269-restoringcoefficientquantity, def:physics:phyx-mini-0269:phyxminiproblems-problemphyxmini0269-nonnegativesireadout}
  Newton-meter readout of a restoring-torque coefficient
\end{definition}
\begin{proof}
  This is the declared model definition of restoring Coefficient In Newton Meters.
\end{proof}
\begin{definition}[angular Velocity In Radians Per Second]
  \label{def:physics:phyx-mini-0269:phyxminiproblems-problemphyxmini0269-angularvelocityinradianspersecond}
  \lean{PhyXMiniProblems.ProblemPhyXMini0269.angularVelocityInRadiansPerSecond}
  \uses{def:physics:phyx-mini-0269:phyxminiproblems-problemphyxmini0269-angularvelocityquantity, def:physics:phyx-mini-0269:phyxminiproblems-problemphyxmini0269-signedsireadout}
  Radian-per-second readout of a signed angular velocity
\end{definition}
\begin{proof}
  This is the declared model definition of angular Velocity In Radians Per Second.
\end{proof}
\begin{definition}[angular Acceleration In Radians Per Second Squared]
  \label{def:physics:phyx-mini-0269:phyxminiproblems-problemphyxmini0269-angularaccelerationinradianspersecondsquared}
  \lean{PhyXMiniProblems.ProblemPhyXMini0269.angularAccelerationInRadiansPerSecondSquared}
  \uses{def:physics:phyx-mini-0269:phyxminiproblems-problemphyxmini0269-angularaccelerationquantity, def:physics:phyx-mini-0269:phyxminiproblems-problemphyxmini0269-signedsireadout}
  Radian-per-second-squared readout of angular acceleration
\end{definition}
\begin{proof}
  This is the declared model definition of angular Acceleration In Radians Per Second Squared.
\end{proof}
\begin{definition}[angular Frequency In Radians Per Second]
  \label{def:physics:phyx-mini-0269:phyxminiproblems-problemphyxmini0269-angularfrequencyinradianspersecond}
  \lean{PhyXMiniProblems.ProblemPhyXMini0269.angularFrequencyInRadiansPerSecond}
  \uses{def:physics:phyx-mini-0269:phyxminiproblems-problemphyxmini0269-angularfrequencyquantity, def:physics:phyx-mini-0269:phyxminiproblems-problemphyxmini0269-nonnegativesireadout}
  Radian-per-second readout of a nonnegative angular frequency
\end{definition}
\begin{proof}
  This is the declared model definition of angular Frequency In Radians Per Second.
\end{proof}
\begin{definition}[angle In Degrees]
  \label{def:physics:phyx-mini-0269:phyxminiproblems-problemphyxmini0269-angleindegrees}
  \lean{PhyXMiniProblems.ProblemPhyXMini0269.angleInDegrees}
  Convert a dimensionless radian coordinate to the degree readout used by the graph
\end{definition}
\begin{proof}
  This is the declared model definition of angle In Degrees.
\end{proof}
\begin{definition}[Plate Material]
  \label{def:physics:phyx-mini-0269:phyxminiproblems-problemphyxmini0269-platematerial}
  \lean{PhyXMiniProblems.ProblemPhyXMini0269.PlateMaterial}
  Material category of the rotating plate
\end{definition}
\begin{proof}
  This is the declared model definition of Plate Material.
\end{proof}
\begin{definition}[Axle Placement]
  \label{def:physics:phyx-mini-0269:phyxminiproblems-problemphyxmini0269-axleplacement}
  \lean{PhyXMiniProblems.ProblemPhyXMini0269.AxlePlacement}
  Position of the plate's axle
\end{definition}
\begin{proof}
  This is the declared model definition of Axle Placement.
\end{proof}
\begin{definition}[Spring Anchor]
  \label{def:physics:phyx-mini-0269:phyxminiproblems-problemphyxmini0269-springanchor}
  \lean{PhyXMiniProblems.ProblemPhyXMini0269.SpringAnchor}
  Body to which the fixed end of the spring is connected
\end{definition}
\begin{proof}
  This is the declared model definition of Spring Anchor.
\end{proof}
\begin{definition}[Spring Attachment]
  \label{def:physics:phyx-mini-0269:phyxminiproblems-problemphyxmini0269-springattachment}
  \lean{PhyXMiniProblems.ProblemPhyXMini0269.SpringAttachment}
  Position of the moving end of the spring on the plate
\end{definition}
\begin{proof}
  This is the declared model definition of Spring Attachment.
\end{proof}
\begin{definition}[Spring Force Model]
  \label{def:physics:phyx-mini-0269:phyxminiproblems-problemphyxmini0269-springforcemodel}
  \lean{PhyXMiniProblems.ProblemPhyXMini0269.SpringForceModel}
  Constitutive model used for the spring
\end{definition}
\begin{proof}
  This is the declared model definition of Spring Force Model.
\end{proof}
\begin{definition}[Angular Motion Regime]
  \label{def:physics:phyx-mini-0269:phyxminiproblems-problemphyxmini0269-angularmotionregime}
  \lean{PhyXMiniProblems.ProblemPhyXMini0269.AngularMotionRegime}
  Motion regime described by the source
\end{definition}
\begin{proof}
  This is the declared model definition of Angular Motion Regime.
\end{proof}
\begin{definition}[Graph Time Unit]
  \label{def:physics:phyx-mini-0269:phyxminiproblems-problemphyxmini0269-graphtimeunit}
  \lean{PhyXMiniProblems.ProblemPhyXMini0269.GraphTimeUnit}
  Units printed on the time axis of the supplied graph
\end{definition}
\begin{proof}
  This is the declared model definition of Graph Time Unit.
\end{proof}
\begin{definition}[Graph Angle Unit]
  \label{def:physics:phyx-mini-0269:phyxminiproblems-problemphyxmini0269-graphangleunit}
  \lean{PhyXMiniProblems.ProblemPhyXMini0269.GraphAngleUnit}
  Units printed on the angular-position axis of the supplied graph
\end{definition}
\begin{proof}
  This is the declared model definition of Graph Angle Unit.
\end{proof}
\begin{definition}[Figure Component]
  \label{def:physics:phyx-mini-0269:phyxminiproblems-problemphyxmini0269-figurecomponent}
  \lean{PhyXMiniProblems.ProblemPhyXMini0269.FigureComponent}
  Individually visible components and labels in the two-panel source image
\end{definition}
\begin{proof}
  This is the declared model definition of Figure Component.
\end{proof}
\begin{definition}[Supplied Figure Readout]
  \label{def:physics:phyx-mini-0269:phyxminiproblems-problemphyxmini0269-suppliedfigurereadout}
  \lean{PhyXMiniProblems.ProblemPhyXMini0269.SuppliedFigureReadout}
  \uses{def:physics:phyx-mini-0269:phyxminiproblems-problemphyxmini0269-lengthquantity, def:physics:phyx-mini-0269:phyxminiproblems-problemphyxmini0269-timequantity, def:physics:phyx-mini-0269:phyxminiproblems-problemphyxmini0269-graphtimeunit, def:physics:phyx-mini-0269:phyxminiproblems-problemphyxmini0269-graphangleunit, def:physics:phyx-mini-0269:phyxminiproblems-problemphyxmini0269-figurecomponent}
  Primary-image data. The right panel begins at a positive angular maximum and returns to the next positive maximum at the grid line labelled t\_s. Thus the named separation is a directly read figure quantity, not a value defined from the requested inertia
\end{definition}
\begin{proof}
  This is the declared model definition of Supplied Figure Readout.
\end{proof}
\begin{definition}[Plate Spring Setup]
  \label{def:physics:phyx-mini-0269:phyxminiproblems-problemphyxmini0269-platespringsetup}
  \lean{PhyXMiniProblems.ProblemPhyXMini0269.PlateSpringSetup}
  \uses{def:physics:phyx-mini-0269:phyxminiproblems-problemphyxmini0269-lengthquantity, def:physics:phyx-mini-0269:phyxminiproblems-problemphyxmini0269-signedlengthquantity, def:physics:phyx-mini-0269:phyxminiproblems-problemphyxmini0269-timequantity, def:physics:phyx-mini-0269:phyxminiproblems-problemphyxmini0269-springconstantquantity, def:physics:phyx-mini-0269:phyxminiproblems-problemphyxmini0269-signedforcequantity, def:physics:phyx-mini-0269:phyxminiproblems-problemphyxmini0269-rotationalinertiaquantity, def:physics:phyx-mini-0269:phyxminiproblems-problemphyxmini0269-signedtorquequantity, def:physics:phyx-mini-0269:phyxminiproblems-problemphyxmini0269-restoringcoefficientquantity, def:physics:phyx-mini-0269:phyxminiproblems-problemphyxmini0269-angularvelocityquantity, def:physics:phyx-mini-0269:phyxminiproblems-problemphyxmini0269-angularaccelerationquantity, def:physics:phyx-mini-0269:phyxminiproblems-problemphyxmini0269-angularfrequencyquantity, def:physics:phyx-mini-0269:phyxminiproblems-problemphyxmini0269-platematerial, def:physics:phyx-mini-0269:phyxminiproblems-problemphyxmini0269-axleplacement, def:physics:phyx-mini-0269:phyxminiproblems-problemphyxmini0269-springanchor, def:physics:phyx-mini-0269:phyxminiproblems-problemphyxmini0269-springattachment, def:physics:phyx-mini-0269:phyxminiproblems-problemphyxmini0269-springforcemodel, def:physics:phyx-mini-0269:phyxminiproblems-problemphyxmini0269-angularmotionregime, def:physics:phyx-mini-0269:phyxminiproblems-problemphyxmini0269-suppliedfigurereadout}
  The independent physical state and response fields. In particular, rotationalInertia, restoringCoefficient, angularFrequency, and motionPeriod are independent quantities constrained only by the laws below; none is defined to equal the requested answer
\end{definition}
\begin{proof}
  This is the declared model definition of Plate Spring Setup.
\end{proof}
\begin{definition}[Is In Small Angle Regime]
  \label{def:physics:phyx-mini-0269:phyxminiproblems-problemphyxmini0269-isinsmallangleregime}
  \lean{PhyXMiniProblems.ProblemPhyXMini0269.IsInSmallAngleRegime}
  \uses{def:physics:phyx-mini-0269:phyxminiproblems-problemphyxmini0269-timequantity, def:physics:phyx-mini-0269:phyxminiproblems-problemphyxmini0269-platespringsetup}
  The interval on which the first-order tangential geometry is assumed
\end{definition}
\begin{proof}
  This is the declared model definition of Is In Small Angle Regime.
\end{proof}
\begin{definition}[Matches Verbal Scenario]
  \label{def:physics:phyx-mini-0269:phyxminiproblems-problemphyxmini0269-matchesverbalscenario}
  \lean{PhyXMiniProblems.ProblemPhyXMini0269.MatchesVerbalScenario}
  \uses{def:physics:phyx-mini-0269:phyxminiproblems-problemphyxmini0269-platespringsetup}
  Qualitative assumptions stated for the apparatus and motion
\end{definition}
\begin{proof}
  This is the declared model definition of Matches Verbal Scenario.
\end{proof}
\begin{definition}[Matches Problem And Figure Readouts]
  \label{def:physics:phyx-mini-0269:phyxminiproblems-problemphyxmini0269-matchesproblemandfigurereadouts}
  \lean{PhyXMiniProblems.ProblemPhyXMini0269.MatchesProblemAndFigureReadouts}
  \uses{def:physics:phyx-mini-0269:phyxminiproblems-problemphyxmini0269-lengthinmeters, def:physics:phyx-mini-0269:phyxminiproblems-problemphyxmini0269-signedlengthinmeters, def:physics:phyx-mini-0269:phyxminiproblems-problemphyxmini0269-timeinseconds, def:physics:phyx-mini-0269:phyxminiproblems-problemphyxmini0269-springconstantinnewtonspermeter, def:physics:phyx-mini-0269:phyxminiproblems-problemphyxmini0269-angularvelocityinradianspersecond, def:physics:phyx-mini-0269:phyxminiproblems-problemphyxmini0269-angleindegrees, def:physics:phyx-mini-0269:phyxminiproblems-problemphyxmini0269-platespringsetup, def:physics:phyx-mini-0269:phyxminiproblems-problemphyxmini0269-isinsmallangleregime}
  Numerical data and readouts supplied by the statement and primary figure: k = 2000 N/m; r = 2.5 cm = 1/40 m; release from 7° and from rest; t\_s = 20 ms = 1/50 s; the graph's next positive maximum occurs one t\_s after the first. The last item is a period measurement from the graph, not an inertia formula
\end{definition}
\begin{proof}
  This is the declared model definition of Matches Problem And Figure Readouts.
\end{proof}
\begin{definition}[Has Physical Plate Spring Parameters]
  \label{def:physics:phyx-mini-0269:phyxminiproblems-problemphyxmini0269-hasphysicalplatespringparameters}
  \lean{PhyXMiniProblems.ProblemPhyXMini0269.HasPhysicalPlateSpringParameters}
  \uses{def:physics:phyx-mini-0269:phyxminiproblems-problemphyxmini0269-lengthinmeters, def:physics:phyx-mini-0269:phyxminiproblems-problemphyxmini0269-timeinseconds, def:physics:phyx-mini-0269:phyxminiproblems-problemphyxmini0269-springconstantinnewtonspermeter, def:physics:phyx-mini-0269:phyxminiproblems-problemphyxmini0269-rotationalinertiainkilogrammeterssquared, def:physics:phyx-mini-0269:phyxminiproblems-problemphyxmini0269-restoringcoefficientinnewtonmeters, def:physics:phyx-mini-0269:phyxminiproblems-problemphyxmini0269-angularfrequencyinradianspersecond, def:physics:phyx-mini-0269:phyxminiproblems-problemphyxmini0269-platespringsetup}
  Strict positivity and nondegeneracy assumptions for the physical model
\end{definition}
\begin{proof}
  This is the declared model definition of Has Physical Plate Spring Parameters.
\end{proof}
\begin{definition}[Satisfies Linearized Plate Spring Laws]
  \label{def:physics:phyx-mini-0269:phyxminiproblems-problemphyxmini0269-satisfieslinearizedplatespringlaws}
  \lean{PhyXMiniProblems.ProblemPhyXMini0269.SatisfiesLinearizedPlateSpringLaws}
  \uses{def:physics:phyx-mini-0269:phyxminiproblems-problemphyxmini0269-lengthinmeters, def:physics:phyx-mini-0269:phyxminiproblems-problemphyxmini0269-signedlengthinmeters, def:physics:phyx-mini-0269:phyxminiproblems-problemphyxmini0269-springconstantinnewtonspermeter, def:physics:phyx-mini-0269:phyxminiproblems-problemphyxmini0269-signedforceinnewtons, def:physics:phyx-mini-0269:phyxminiproblems-problemphyxmini0269-rotationalinertiainkilogrammeterssquared, def:physics:phyx-mini-0269:phyxminiproblems-problemphyxmini0269-signedtorqueinnewtonmeters, def:physics:phyx-mini-0269:phyxminiproblems-problemphyxmini0269-restoringcoefficientinnewtonmeters, def:physics:phyx-mini-0269:phyxminiproblems-problemphyxmini0269-angularaccelerationinradianspersecondsquared, def:physics:phyx-mini-0269:phyxminiproblems-problemphyxmini0269-platespringsetup, def:physics:phyx-mini-0269:phyxminiproblems-problemphyxmini0269-isinsmallangleregime}
  First-order mechanics for the small angular motion. The positive angular direction is chosen to increase the spring length: tangential extension is x = r θ; Hooke's law is F = -k x; the perpendicular rim lever arm gives τ = r F; the spring supplies the total axle torque; τ = -κ θ defines the effective restoring coefficient; rotational Newton's law is I α = τ. These are governing laws. No field solves for I, fixes it numerically, or selects an answer choice
\end{definition}
\begin{proof}
  This is the declared model definition of Satisfies Linearized Plate Spring Laws.
\end{proof}
\begin{definition}[scalar Rotational Oscillator]
  \label{def:physics:phyx-mini-0269:phyxminiproblems-problemphyxmini0269-scalarrotationaloscillator}
  \lean{PhyXMiniProblems.ProblemPhyXMini0269.scalarRotationalOscillator}
  \uses{def:physics:phyx-mini-0269:phyxminiproblems-problemphyxmini0269-rotationalinertiainkilogrammeterssquared, def:physics:phyx-mini-0269:phyxminiproblems-problemphyxmini0269-restoringcoefficientinnewtonmeters, def:physics:phyx-mini-0269:phyxminiproblems-problemphyxmini0269-platespringsetup, def:physics:phyx-mini-0269:phyxminiproblems-problemphyxmini0269-hasphysicalplatespringparameters}
  The scalar equation I θ'' + κ θ = 0 has the form represented by Physlib's positive ClassicalMechanics.HarmonicOscillator. The oscillator's m slot is therefore filled by the SI readout of rotational inertia, and its k slot by the SI readout of the restoring-torque coefficient
\end{definition}
\begin{proof}
  This is the declared model definition of scalar Rotational Oscillator.
\end{proof}
\begin{definition}[Satisfies Ideal Angular SHMLaw]
  \label{def:physics:phyx-mini-0269:phyxminiproblems-problemphyxmini0269-satisfiesidealangularshmlaw}
  \lean{PhyXMiniProblems.ProblemPhyXMini0269.SatisfiesIdealAngularSHMLaw}
  \uses{def:physics:phyx-mini-0269:phyxminiproblems-problemphyxmini0269-timeinseconds, def:physics:phyx-mini-0269:phyxminiproblems-problemphyxmini0269-angularfrequencyinradianspersecond, def:physics:phyx-mini-0269:phyxminiproblems-problemphyxmini0269-platespringsetup, def:physics:phyx-mini-0269:phyxminiproblems-problemphyxmini0269-hasphysicalplatespringparameters, def:physics:phyx-mini-0269:phyxminiproblems-problemphyxmini0269-scalarrotationaloscillator}
  The measured frequency and period agree with Physlib's generic harmonic oscillator, and one period is the separation of consecutive positive peaks in the graph. This predicate contains no expansion of the unknown inertia in terms of this problem's numerical data
\end{definition}
\begin{proof}
  This is the declared model definition of Satisfies Ideal Angular SHMLaw.
\end{proof}
\begin{lemma}[restoring Coefficient formula]
  \label{lem:physics:phyx-mini-0269:phyxminiproblems-problemphyxmini0269-restoringcoefficient-formula}
  \lean{PhyXMiniProblems.ProblemPhyXMini0269.restoringCoefficient_formula}
  \uses{def:physics:phyx-mini-0269:phyxminiproblems-problemphyxmini0269-lengthinmeters, def:physics:phyx-mini-0269:phyxminiproblems-problemphyxmini0269-springconstantinnewtonspermeter, def:physics:phyx-mini-0269:phyxminiproblems-problemphyxmini0269-restoringcoefficientinnewtonmeters, def:physics:phyx-mini-0269:phyxminiproblems-problemphyxmini0269-platespringsetup, def:physics:phyx-mini-0269:phyxminiproblems-problemphyxmini0269-matchesproblemandfigurereadouts, def:physics:phyx-mini-0269:phyxminiproblems-problemphyxmini0269-satisfieslinearizedplatespringlaws}
  The figure geometry, small-angle extension, Hooke force, and lever-arm torque imply the intermediate relation κ = k r². It remains a conclusion rather than a law field because it combines several independently stated laws
\end{lemma}
\begin{proof}
  The conclusion follows from the governing relations and figure readouts named in the statement of restoring Coefficient formula.
\end{proof}
\begin{lemma}[rotational Inertia exact]
  \label{lem:physics:phyx-mini-0269:phyxminiproblems-problemphyxmini0269-rotationalinertia-exact}
  \lean{PhyXMiniProblems.ProblemPhyXMini0269.rotationalInertia_exact}
  \uses{def:physics:phyx-mini-0269:phyxminiproblems-problemphyxmini0269-lengthinmeters, def:physics:phyx-mini-0269:phyxminiproblems-problemphyxmini0269-timeinseconds, def:physics:phyx-mini-0269:phyxminiproblems-problemphyxmini0269-springconstantinnewtonspermeter, def:physics:phyx-mini-0269:phyxminiproblems-problemphyxmini0269-rotationalinertiainkilogrammeterssquared, def:physics:phyx-mini-0269:phyxminiproblems-problemphyxmini0269-platespringsetup, def:physics:phyx-mini-0269:phyxminiproblems-problemphyxmini0269-matchesproblemandfigurereadouts, def:physics:phyx-mini-0269:phyxminiproblems-problemphyxmini0269-hasphysicalplatespringparameters, def:physics:phyx-mini-0269:phyxminiproblems-problemphyxmini0269-satisfieslinearizedplatespringlaws, def:physics:phyx-mini-0269:phyxminiproblems-problemphyxmini0269-satisfiesidealangularshmlaw}
  For κ = k r² and T = 2π sqrt(I/κ), the rotational inertia is I = k r² T² / (4π²). This is the exact physical relation inferred from the spring and graph; no hypothesis above contains it
\end{lemma}
\begin{proof}
  The conclusion follows from the governing relations and figure readouts named in the statement of rotational Inertia exact.
\end{proof}
\begin{definition}[Answer Choice]
  \label{def:physics:phyx-mini-0269:phyxminiproblems-problemphyxmini0269-answerchoice}
  \lean{PhyXMiniProblems.ProblemPhyXMini0269.AnswerChoice}
  Labels of the four rotational-inertia choices printed with the problem
\end{definition}
\begin{proof}
  This is the declared model definition of Answer Choice.
\end{proof}
\begin{definition}[displayed Inertia In Kilogram Meters Squared]
  \label{def:physics:phyx-mini-0269:phyxminiproblems-problemphyxmini0269-displayedinertiainkilogrammeterssquared}
  \lean{PhyXMiniProblems.ProblemPhyXMini0269.displayedInertiaInKilogramMetersSquared}
  \uses{def:physics:phyx-mini-0269:phyxminiproblems-problemphyxmini0269-answerchoice}
  Displayed rotational inertia for each choice, in kg · m²
\end{definition}
\begin{proof}
  This is the declared model definition of displayed Inertia In Kilogram Meters Squared.
\end{proof}
\begin{definition}[recorded Dataset Answer]
  \label{def:physics:phyx-mini-0269:phyxminiproblems-problemphyxmini0269-recordeddatasetanswer}
  \lean{PhyXMiniProblems.ProblemPhyXMini0269.recordedDatasetAnswer}
  \uses{def:physics:phyx-mini-0269:phyxminiproblems-problemphyxmini0269-answerchoice}
  The dataset's recorded answer label; metadata, not a theorem premise
\end{definition}
\begin{proof}
  This is the declared model definition of recorded Dataset Answer.
\end{proof}
\begin{definition}[Is Unique Closest Answer]
  \label{def:physics:phyx-mini-0269:phyxminiproblems-problemphyxmini0269-isuniqueclosestanswer}
  \lean{PhyXMiniProblems.ProblemPhyXMini0269.IsUniqueClosestAnswer}
  \uses{def:physics:phyx-mini-0269:phyxminiproblems-problemphyxmini0269-rotationalinertiainkilogrammeterssquared, def:physics:phyx-mini-0269:phyxminiproblems-problemphyxmini0269-platespringsetup, def:physics:phyx-mini-0269:phyxminiproblems-problemphyxmini0269-answerchoice, def:physics:phyx-mini-0269:phyxminiproblems-problemphyxmini0269-displayedinertiainkilogrammeterssquared}
  A listed choice is uniquely closest to the exact inferred SI value. This models the rounding implicit in a multiple-choice answer instead of falsely asserting that the irrational exact value is equal to a displayed decimal
\end{definition}
\begin{proof}
  This is the declared model definition of Is Unique Closest Answer.
\end{proof}
% --- Archon declaration topology end ---
% --- Archon physics formalization source end ---
